\honest{An Especially Elegant Evpsych Experiment}{Eliezer Yudkowsky}{2009}

I quote Robert Wright's summary of a 1989 Canadian study. Adults judged which deaths of children of various ages would most grieve a parent. Their grief curve correlated .64 with the reproductive value of children at each age in Canada, and ``an extremely high .92'' with the same curve for the !Kung, a hunter-gatherer people. \nb{Wright's summary checks against the book it comes from, \textsc{The Moral Animal}. The paper's third author is K. L. Jang; the post has ``Lang.''} At first I have not read the paper; then a reader finds it, and the correlations prove to be as described, ``N=221.''

The study's weakness, I grant, is that it asked adults to imagine a parent's grief instead of asking parents. ``My understanding is'' that closer studies agree. \nb{The post names none. A 2020 review by Reynolds and colleagues cites two: in one, bereaved parents grieved most for a child lost at 17; in the other, bereaved mothers grieved most for adolescents.}

Then the merits. Selection strong enough to create grief could tune it this finely. The study answers people who say evolutionary psychology makes no advance predictions: no one would think of running it without evolutionary psychology. And it shows the difference between a hidden motive and an adaptation that simply runs: grief is ``not even subconsciously'' about reproductive value, or it would follow the Canadian curve. ``Parents care about children for their own sake''; that children carry their genes is the evolutionary, not the mental, reason such minds exist. I draw my cynicism about hidden motives sharply at that line, or I might as well tell the supermarket cashier that bagging my groceries correctly is only a way to maximize inclusive fitness. \nb{The study compared imagined grief with two demographic curves; it did not measure what parents care about. Robin Hanson objected in the comments that rejecting one extreme ``does not imply the other pole,'' and Yudkowsky replied, ``I wasn't arguing for the other pole.''}

.92 is, I think, among the highest correlations in psychology; the only higher one I recall, .98, came from two groups answering what amounts to the same question, similarity and probability, and ``correlation .92 and N=221 is pretty strong evidence.'' \nb{The 221 are raters. The correlation is between two curves, one point per age; a commenter who read the paper counted ten ages. On ten points .92 is still unlikely by chance, but a correlation between averages over hundreds of raters is not comparable with correlations between individuals. The post also does not say whether the gap between .92 and .64 could be chance.} The fine-tuning I take for granted, given the huge selection pressure. The remaining doubts are whether adults can picture relative grief correctly, and whether the !Kung are typical of ancestral hunter-gatherers or tribes varied too much for .92.

Imagined grief is the right thing to measure, because it is what steers a parent before a loss; to evolution a dead child is a sunk cost, and it ``wants'' the parent to learn, recover and raise other children. \nb{The subjects were not imagining their own loss; they estimated what a parent would feel.} The curve that fits is the child's future reproductive value, not what has already been spent on the child; perhaps the fit would be even better with the future cost of raising the child to independence. Humans notice sunk costs, presumably to keep us from switching strategies too often, or as an unfortunate spandrel of the pain of waste. Natural selection does not; it does not ``think'' at all, being only a fact about past reproduction. Yet the grief adaptation goes on working as if the parent lived among the !Kung, not in Canada, a difference most humans would notice. ``Humans and natural selection are insane in different stable complicated ways.''
